\chapter{Diskussion}\label{chap:Diskussion}

Die vier Modelle liegen im gewählten Testsetup eng beieinander. Das MLP erreicht den höchsten Macro-F1-Wert von 0{,}934. HistGradientBoosting folgt mit 0{,}932, Random Forest und logistische Regression mit 0{,}931. Der Vorsprung des MLP beträgt damit nur 0{,}003 gegenüber dem Random Forest. Die Auswahl eines Modells sollte daher nicht allein an der dritten Nachkommastelle festgemacht werden.

Die XAI-Ergebnisse ergänzen diesen Befund. Über alle Modelle hinweg stehen \texttt{Area} und \texttt{ShapeFactor1} häufig weit oben. Weitere wichtige Merkmale hängen vom Modell ab. Random Forest und HistGradientBoosting heben stärker Rundheit und Kompaktheit hervor. Beim MLP und bei der logistischen Regression spielt \texttt{AspectRatio} in der globalen Erklärung eine größere Rolle. Der Befund passt zur Aufgabenstellung, weil die ausgewählten Variablen Größe und geometrische Form der Bohnen abbilden. Er beweist aber nicht, dass genau diese Merkmale die botanischen Ursachen der Sortenzugehörigkeit sind.

Die Unterschiede zwischen PFI und SHAP sind erwartbar. PFI fragt nach dem Verlust des Macro-F1 nach einer Störung. SHAP beschreibt Beiträge zu einzelnen Modellvorhersagen und verdichtet sie danach. Bei korrelierten Merkmalen ist keine dieser Rangfolgen eine eindeutige Zuordnung von Information zu einer einzelnen Variable. Molnar weist besonders für PFI auf problematische, durch Permutation erzeugte Merkmalskombinationen hin \cite[Kap.~23]{molnar2025interpretable}. Die XAI-Befunde werden deshalb als wiederkehrendes Muster gelesen, nicht als exakte Reihenfolge.

Auch die lokalen Beispiele haben eine Grenze. Die korrekt klassifizierte DERMASON-Instanz und die DERMASON-als-SIRA-Fehler machen Entscheidungen greifbar, repräsentieren aber nicht die gesamte Klasse. Dass alle Modelle einen solchen Fehlfall besitzen, bestätigt die Schwierigkeit der Trennung. Die Konfusionsmatrizen bleiben für die Häufigkeit von Fehlern die passendere globale Sicht.

Die Untersuchung verwendet nur acht aus 16 verfügbaren Merkmalen und einen festen stratifizierten Train-Test-Split. Die Ergebnisse gelten deshalb für dieses Feature-Set und diese Aufteilung. Eine externe Validierung oder wiederholte Cross-Validation wurde nicht durchgeführt. Außerdem sind die Eingangsvariablen bereits aus Bildauswertung abgeleitet. Aussagen betreffen folglich die Klassifikation auf diesen morphologischen Messwerten, nicht eine vollständige Bilderkennung unter wechselnden Aufnahmebedingungen.
